\subsection{Numerical restrictions at the smallest endpoint}
\label{subsec:smallest-endpoint-numerical}

The nontrivial discrepancy line at logarithmic order two restricts
the space of cubic sections of the normal bundle.  The hyperplane-section
estimate then gives a lower bound for the degree of a very ample
endpoint.  A second restriction arises when the Albanese variety is a
surface and its Albanese morphism has degree two: the normalized torsor
is then birational to an abelian surface.

Throughout this subsection, let $k$ be algebraically closed of
characteristic two, let $X$ be a smooth integral projective surface,
and let $C\subset X$ be a smooth integral ample divisor.  Put
\[
 d=C^2,\qquad N=\OO_C(C).
\]
Let $M\ne\OO_X$ have exact order two, be trivial on $C$, and admit a
square root; equivalently, $[M]\in2\operatorname{NS}(X)$.
Assume that its logarithmic form has exact order two along $C$ and
that
\[
 2g(C)-2=3d.
\]
We use the notation of
Theorem~\ref{thm:smallest-endpoint-global-structure}:
$\pi:V\to X$ is the normalized torsor, $B=\pi^*C$, and
$\mathcal T=\omega_C\otimes N^{-3}$.
In particular, $d$ is even, $\mathcal T$ has exact order two, and
$K_X\cdot C=2d$.

\begin{proposition}[The degree of a very ample endpoint]
\label{prop:smallest-endpoint-very-ample-degree}
If $C$ is very ample, then
\[
 d\ge12,\qquad
 h^0(C,N)\le\left\lfloor\frac d4\right\rfloor+1,\qquad
 h^0(X,\OO_X(C))\le\left\lfloor\frac d4\right\rfloor+2.
\]
More generally, an embedding of $C$ by a subspace of $H^0(C,N)$
as a nondegenerate curve in $\PP^n$, with $n\ge2$, satisfies
$d\ge4n$.
\end{proposition}

\begin{proof}
Since $\mathcal T$ is a nontrivial degree-zero line bundle, Serre
duality and Riemann--Roch give
\begin{equation}\label{eq:smallest-endpoint-cubic-dimension}
 h^1(C,N^3)=h^0(C,\mathcal T)=0,\qquad
 h^0(C,N^3)=3d+1-g(C)=\frac{3d}{2}.
\end{equation}

We apply the hyperplane-section argument in the proof of
Lemma~\ref{lem:small-degree-curve-bound} before its final
small-degree specialization.  That argument applies to every
$n\ge2$.  Indeed, the twisted tangent--normal sequence shows that a
smooth embedded curve of genus at least two is nonstrange: a common
tangent point would give a nonzero section of $T_C\otimes N^{-1}$,
whose degree is negative.  Kadets's sectional-monodromy theorem
holds for $n\ge2$
\cite[Proposition~2.8 and proof, pp.~4--5]{Kadets}.
Thus the geometric generic hyperplane section $\Gamma$ consists
of $d$ points in linear general position in $\PP^{n-1}$.
The spanning and interpolation arguments in the cited proof give
\[
 d\ge n,\qquad
 h_\Gamma(j)\ge\min\{d,j(n-1)+1\}\quad(j\ge1).
\]
If $h_C$ denotes the homogeneous-coordinate Hilbert function,
multiplication by a hyperplane equation is injective in the
coordinate ring of $C$, and its quotient surjects onto the
coordinate ring of $\Gamma$.  Consequently
\begin{equation}\label{eq:smallest-endpoint-cubic-bound}
 \begin{split}
 \frac{3d}{2}
 &=h^0(C,N^3)\ge h_C(3)\\
 &\ge 1+n+\min\{d,2n-1\}+\min\{d,3n-2\}.
 \end{split}
\end{equation}
The inequality uses neither projective normality nor uniform
position in higher degree.

If $d<2n-1$, the last expression is $1+n+2d$, which is greater
than $3d/2$.  If $2n-1\le d<3n-2$, it is $3n+d$, again greater
than $3d/2$.  Hence $d\ge3n-2$, and
\eqref{eq:smallest-endpoint-cubic-bound} gives
$3d/2\ge6n-2$.  Since $d$ is even, this implies $d\ge4n$.

The restriction $N$ of the very ample bundle $\OO_X(C)$ is very
ample.  Its complete embedding cannot have image in $\PP^1$,
because $g(C)\ge2$.  If its image were a plane curve, adjunction
would give $d(d-3)=3d$, hence $d=6$ and $\omega_C=N^3$.
This contradicts $\mathcal T\ne\OO_C$.
Thus $n=h^0(C,N)-1\ge3$, proving the degree and normal-section
bounds.  Finally the kernel of
\[
 H^0(X,\OO_X(C))\longrightarrow H^0(C,N)
\]
is the one-dimensional space generated by the defining section
of $C$.  This proves the bound on sections on $X$.
\end{proof}

\begin{proposition}[An Albanese morphism of degree two]
\label{prop:smallest-endpoint-albanese-degree-two}
Suppose in addition that $M$ is algebraically trivial,
$A=\operatorname{Alb}(X)$ has dimension two, and the Albanese
morphism $\alpha:X\to A$ is generically finite of degree two.
Then $V$ is birational to the abelian surface
$A_V=\operatorname{Alb}(V)$, the canonical divisor $K_X$ is nef
and big, and
\begin{equation}\label{eq:smallest-endpoint-albanese-chern-interval}
 d=4\chi(\OO_X),\qquad
 2d\le K_X^2<3d,\qquad
 0<c_2(X)=3d-K_X^2\le d.
\end{equation}
In particular, $4\mid d$ and $X$ is a minimal surface of general type.
\end{proposition}

\begin{proof}
Let $\beta:V\to A_V$ be the Albanese morphism.
Proposition~\ref{prop:smallest-endpoint-albanese-quotient}
identifies $A_V^{(2)}$ with $A$ and gives
\[
 \alpha\pi=F_{A_V/k}\beta.
\]
Relative Frobenius of the abelian surface has degree four, whereas
$\pi$ has degree two.  Thus
\[
 2\deg\alpha=4\deg\beta.
\]
It follows that $\deg\beta=1$, so $\beta$ is proper and birational.

We first compute $\chi(\OO_V)$ without assuming that $V$ is smooth.
Choose a projective resolution $r:S\to V$ and put $b=\beta r$.
The proper birational morphism $b$ between smooth surfaces is a
composition of blowups at closed points
\cite[Lemma~54.17.1, Tag~0C5Q]{Stacks}.
A point blowup has direct image $\OO$ and vanishing higher direct
images of $\OO$.  For the first higher direct image, this follows
on an affine neighborhood with regular parameters $x,y$ from the
two affine blowup charts: every Laurent expression in $y/x$ on
their overlap is a sum of expressions on the two charts.
The resulting first \v{C}ech cohomology vanishes; higher direct
images above degree one vanish by fibre dimension.
These statements persist under composition.
Normality gives $r_*\OO_S=\OO_V$ and
$\beta_*\OO_V=\OO_{A_V}$.
The low-degree Leray sequence gives
\[
 R^1\beta_*\OO_V\hookrightarrow R^1b_*\OO_S=0.
\]
The higher direct images of $\beta$ vanish above degree one.
Consequently
\[
 \chi(\OO_V)=\chi(\OO_{A_V})=0.
\]
The Euler-characteristic formula in
Theorem~\ref{thm:smallest-endpoint-global-structure} now gives
$d=4\chi(\OO_X)$.

Pull back a nowhere-zero two-form on $A_V$ to the regular locus
of $V$.  It extends over the omitted codimension-two set because
$\omega_V$ is reflexive.  Its zero divisor $E_V$ is effective
and supported on the exceptional locus of $\beta$.
It is nonzero, since
\[
 E_V\cdot B=K_V\cdot B
 =2(K_X-C)\cdot C=2d>0.
\]
Here $V$ is Gorenstein and
$\omega_V\simeq\pi^*(\omega_X(-C))$ by
\eqref{eq:smallest-endpoint-global-conductor}.
On $S$, the Cartier pullback $r^*E_V$ is a nonzero effective
$b$-exceptional divisor.  It is orthogonal to $b^*H$ for an ample
divisor $H$ on $A_V$, and $(b^*H)^2>0$.
The Hodge index theorem gives $(r^*E_V)^2<0$: equality would make
this effective nonzero divisor numerically zero, contradicting
its positive intersection with an ample divisor on $S$.
Therefore
\[
 0>K_V^2=2(K_X-C)^2=2K_X^2-6d.
\]

The finite-flat norm of the nonzero canonical section on $V$
is a nonzero section of $\omega_X^{\otimes2}(-2C)$.
Indeed, the determinant of multiplication defines the norm of
a section, and $\operatorname{Nm}_\pi(\pi^*L)=L^{\otimes2}$
for every line bundle $L$ on $X$.
Thus an effective divisor $E_X$ satisfies
\begin{equation}\label{eq:smallest-endpoint-albanese-canonical-norm}
 2K_X\sim2C+E_X.
\end{equation}
Since $C$ is ample, $K_X$ is big.

To prove nefness, suppose an integral curve $E$ has $K_X\cdot E<0$.
Equation~\eqref{eq:smallest-endpoint-albanese-canonical-norm},
ampleness of $C$, and nonnegative intersection of distinct effective
curves force $E$ to be a component of $E_X$ with $E^2<0$.
Adjunction gives
\[
 2p_a(E)-2=K_X\cdot E+E^2.
\]
Both terms on the right are negative integers, so $p_a(E)\ge0$
forces $K_X\cdot E=E^2=-1$ and $p_a(E)=0$.
Then $E$ is a smooth rational $(-1)$-curve, contrary to
Theorem~\ref{thm:smallest-endpoint-global-structure}.
Thus $K_X$ is nef.  Intersecting
\eqref{eq:smallest-endpoint-albanese-canonical-norm} with $K_X$
now gives $2K_X^2\ge2K_X\cdot C=4d$.
Finally Noether's formula, obtained from
Grothendieck--Riemann--Roch in every characteristic
\cite[Tag~02UO]{Stacks}, gives
\[
 c_2(X)=12\chi(\OO_X)-K_X^2=3d-K_X^2.
\]
This proves \eqref{eq:smallest-endpoint-albanese-chern-interval}.
Nefness and bigness give the general-type assertion, and the absence
of $(-1)$-curves gives minimality.
\end{proof}

\begin{corollary}[Very ample degree in the Albanese degree-two case]
\label{cor:smallest-endpoint-albanese-very-ample-degree}
Under Proposition~\ref{prop:smallest-endpoint-albanese-degree-two},
suppose that $C$ is very ample and let
$r=h^0(X,\OO_X(C))-1$ be the dimension of its complete embedding.
Then
\[
 r\ge5,\qquad 4\mid d,\qquad d\ge16,\qquad
 6\le h^0(X,\OO_X(C))\le d/4+2.
\]
If $d=16$, then
\[
 h^0(X,\OO_X(C))=6,\qquad h^0(C,N)=5.
\]
\end{corollary}

\begin{proof}
An embedded surface has $r\ge2$.  For $r=2$ its image is $\PP^2$,
contrary to $K_X\cdot C=2d>0$.
For $r=3$ it is a degree-$d$ hypersurface, so
$K_X=(d-4)C$.  The same intersection identity gives $d=6$
and $K_X=2C$, contrary to $\mathcal T\ne\OO_C$.

Suppose $r=4$ and write $E=N_{X/\PP^4}$.
The tangent--normal and Euler sequences, followed by the
self-intersection formula, give
\[
 \deg c_2(E)=10d+5K_X\cdot C+K_X^2-c_2(X)=d^2.
\]
These are the identities used in the proof of
Proposition~\ref{prop:very-ample-numerical-obstruction};
the self-intersection statement is valid in every characteristic
\cite[Tags~0FV9 and~0FEZ]{Stacks}.
Using $K_X\cdot C=2d$ and
\eqref{eq:smallest-endpoint-albanese-chern-interval} gives
\begin{equation}\label{eq:smallest-endpoint-p4-double-point}
 d^2=17d+2K_X^2.
\end{equation}
The interval $2d\le K_X^2<3d$ yields $21\le d<23$.
There is no multiple of four in this interval, so $r\ne4$.

Thus $r\ge5$.  Proposition~\ref{prop:smallest-endpoint-very-ample-degree}
gives $r+1\le d/4+2$, hence $d\ge4(r-1)\ge16$.
At $d=16$ the upper and lower section bounds give
$h^0(X,\OO_X(C))=6$.
Restriction to $C$ then gives $h^0(C,N)\ge5$, whereas the same
proposition gives $h^0(C,N)\le5$.
\end{proof}

\begin{remark}[Scope of the numerical bounds]
\label{rem:smallest-endpoint-numerical-scope}
The very ample degree bound has no Albanese hypothesis.
The Chern interval requires both $\dim\operatorname{Alb}(X)=2$
and Albanese degree two.  A surface image in an Albanese variety
of larger dimension, or a larger degree over an abelian surface,
does not give the birational morphism used in its proof.
For a merely ample endpoint the cubic-section bound does not apply.
The degree-sixteen conclusion is a necessary condition and makes
no existence assertion.
\end{remark}
